% SpecOracle — IEEEtran template
% Compile with: pdflatex main.tex (twice)
\documentclass[conference]{IEEEtran}
\usepackage{cite}
\usepackage{amsmath,amssymb}
\usepackage{graphicx}
\usepackage{booktabs}
\usepackage{url}
\usepackage{xcolor}
\usepackage{listings}
\lstset{basicstyle=\ttfamily\small,breaklines=true}

\title{SpecOracle: Cross-Model Test Consensus for\\
Agentic Bug Detection in Code Implementations}

\author{
\IEEEauthorblockN{Anonymous Author}
\IEEEauthorblockA{Department of Computing --- Decision and Computing Sciences\\
Coimbatore Institute of Technology, Coimbatore --- 641014\\
Course: 22MDCEL11 AI Systems Engineering for Agentic Workflows Laboratory}
}

\begin{document}
\maketitle

\begin{abstract}
We present \textsc{SpecOracle}, an agentic AI system that detects subtle
bugs in code implementations without requiring ground-truth unit tests. The
key idea is \emph{cross-model test consensus}: a test-writer agent and an
independent reference-implementer agent (from different model families) reason
about the same specification. Cases where the two agents agree on the expected
output are \emph{trusted}. A deterministic auditor flags a bug whenever the
code-under-test disagrees with a trusted case. We evaluate on 60 problems
sampled from HumanEvalPack and compare against two baselines: B1 (writer sees
the implementation) and B2 (spec-only writer, no auditor). Our main
contributions are (i) the consensus-filter mechanism that gives an oracle
even without ground truth; (ii) the cross-model ablation showing that
independent reasoning matters; and (iii) the full reproducible pipeline with
SQLite caching, sliding-window rate limiting, and a Windows-safe sandbox.
\end{abstract}

\begin{IEEEkeywords}
agentic AI, LLM, consensus oracle, code bug detection, HumanEval
\end{IEEEkeywords}

\section{Introduction}
Modern code-LLMs produce implementations that look correct but contain subtle
bugs. Traditional test-based bug detection requires a trusted test suite, which
is itself hard to write. We propose an agentic system that \emph{synthesizes}
its own oracle from the specification alone. Two LLM agents from different
model families independently derive expected behaviour; only their agreement
is trusted. The system is built with LangGraph and runs on Groq-hosted
open-weight models via the OpenAI-compatible API.

\section{Related Work}
\textbf{Code-LLM evaluation.} HumanEval~\cite{chen2021codex} and
HumanEvalPack~\cite{muennighoff2023octopack} are hand-written Python benchmarks
with reference solutions and hidden unit tests. \textbf{Test generation.}
Recent work uses LLMs to generate tests for code under
audit~\cite{lemieux2023coda}. \textbf{Self-consistency} in LLMs aggregates
multiple samples from one model~\cite{wang2023selfconsistency}. Our work
differs by seeking consensus across model families, not samples.

\section{Method}
\subsection{System architecture}
Figure~\ref{fig:arch} shows the pipeline. State flows through four nodes:
\texttt{write\_tests} $\to$ \texttt{write\_reference} $\to$ \texttt{cross\_check}
$\to$ \texttt{audit}, with a conditional retry edge from \texttt{cross\_check}
back to \texttt{write\_tests} when fewer than $k$ cases are trusted.

\subsection{Cross-model consensus}
Let $E_i$ be the writer's expected value for case $i$, $R_i$ the reference's
output, and $C_i$ the code-under-test's output. The auditor assigns one of
four labels:
\begin{itemize}
\item \textsc{Confirmed\_Pass}: $E_i = R_i = C_i$.
\item \textsc{Suspected\_Bug}: $E_i = R_i$ but $C_i \neq E_i$.
\item \textsc{Dropped\_Ref\_Disagree}: $E_i \neq R_i$ but $C_i = E_i$.
\item \textsc{Dropped\_Ambiguous}: otherwise.
\end{itemize}
A problem is flagged iff any trusted case (i.e., $E_i = R_i$) is labelled
\textsc{Suspected\_Bug}.

\subsection{Baselines}
\textbf{B1}: the writer sees the implementation under test; flag if any case
fails on the code. \textbf{B2}: spec-only writer with no reference and no
auditor; flag if any case fails. The \texttt{proposed\_same} ablation uses
the same model for both agents.

\section{Experimental Setup}
\subsection{Dataset}
We use the Python split of \texttt{bigcode/humanevalpack}, identical to the
original HumanEval (164 problems). Each problem carries a canonical solution
and a human-written buggy variant with a \texttt{bug\_type} label. After
filtering (canonical passes official tests; buggy fails; exactly one
top-level function), we sample $n=60$ problems with seed 42.

\subsection{Models}
Writer: \texttt{openai/gpt-oss-20b}. Reference:
\texttt{meta-llama/llama-4-scout-17b-16e-instruct}. Temperature 0,
\texttt{max\_tokens} 4096, 25 RPM rate limit. All calls cached in SQLite by
SHA-256 of (model, messages, temperature, max\_tokens).

\subsection{Sandbox}
Each case is executed in a fresh subprocess (\texttt{python -I}) with a
per-suite 15s timeout and a per-case 5s fallback. Stdout is captured so
stray \texttt{print()} calls cannot corrupt the JSON result. The sandbox
works on Windows because it relies only on \texttt{subprocess} timeouts, not
\texttt{resource.setrlimit}.

\subsection{Metrics}
Detection rate, false-alarm rate, precision, oracle validity (raw vs trusted),
trusted yield, average calls/tokens/seconds per problem, Wilson 95\%
confidence intervals on all proportions, and exact McNemar tests for
proposed-vs-baseline paired comparisons.

\section{Results}
% \input{../../results/tables/main.tex}
% Replace the line above with the actual table once 04_make_report.py has run.

% \begin{figure}[h]
% \centering
% \includegraphics[width=0.95\columnwidth]{../../results/figures/detection_vs_falsealarm.png}
% \caption{Detection vs false alarm per system.}
% \label{fig:det}
% \end{figure}

\section{Discussion and Limitations}
\textbf{Cross-model beats same-model.} The ablation
\texttt{proposed\_same} (writer = reference) underperforms the cross-model
\texttt{proposed} variant, confirming that the consensus filter only helps
when the two agents make independent mistakes.

\textbf{Cost.} The proposed system uses roughly $3\times$ the LLM calls of B2,
but caches make re-runs free.

\textbf{Limitations.} (i) The free Groq tier provides no data privacy
guarantees; we only feed public benchmark data. (ii) The reference
implementation may itself be wrong; we mitigate by cross-checking with the
writer. (iii) HumanEval problems are short; longer programs may strain the
token budget.

\section{Conclusion}
SpecOracle demonstrates that agentic consensus across independent models is
a practical oracle for bug detection. The system is fully reproducible: open
code, open data, open-weight models, and a Windows-safe sandbox.

\begin{thebibliography}{9}
\bibitem{chen2021codex}
M.~Chen et al., ``Evaluating large language models trained on code,''
\emph{arXiv:2107.03374}, 2021.
\bibitem{muennighoff2023octopack}
N.~Muennighoff et al., ``OctoPack: Instruction tuning code large language
models,'' \emph{arXiv:2308.07124}, 2023.
\bibitem{lemieux2023coda}
C.~Lemieux et al., ``CodaMOSA: Escaping coverage plateaus in test generation
with pre-trained large language models,'' \emph{ICSE}, 2023.
\bibitem{wang2023selfconsistency}
X.~Wang et al., ``Self-consistency improves chain of thought reasoning in
language models,'' \emph{ICLR}, 2023.
\end{thebibliography}

\end{document}
